\section{Segunda Ley de Newton}

La fuerza neta aplicada a un cuerpo es igual a la tasa de cambio de su cantidad de movimiento. Para masa constante:
\begin{equation}
    \sum \vec{F} = m \cdot \vec{a}
\end{equation}

\subsection{Componentes}
\begin{align}
    \sum F_x &= m \cdot a_x \\
    \sum F_y &= m \cdot a_y
\end{align}

\subsection{Diagrama de cuerpo libre (DCL)}
Es una herramienta gráfica que muestra todas las fuerzas que actúan sobre un cuerpo. Pasos para dibujarlo:
\begin{enumerate}
    \item Aislar el cuerpo.
    \item Dibujar todas las fuerzas externas (peso, normal, tensión, fricción, etc.).
    \item Elegir un sistema de coordenadas.
    \item Descomponer las fuerzas en sus componentes.
\end{enumerate}

\section{Tercera Ley de Newton (Acción y Reacción)}

\subsection{Enunciado}
Si un cuerpo A ejerce una fuerza sobre un cuerpo B, entonces B ejerce una fuerza sobre A de igual magnitud, misma dirección y sentido opuesto.
\begin{equation}
    \vec{F}_{AB} = -\vec{F}_{BA}
\end{equation}

\subsection{Implicaciones}
\begin{itemize}
    \item Las fuerzas siempre actúan en pares.
    \item Las fuerzas de acción y reacción nunca se anulan porque actúan sobre cuerpos diferentes.
\end{itemize}

\resumebox{ejercicio}{Fuerza magnética del protón}{\footnotesize
    Un bloque de \(10 \, \text{kg}\) está sobre una superficie horizontal. Se aplica una fuerza de \(50 \, \text{N}\) horizontal. Si el coeficiente de fricción cinética es \(0.2\), calcula la aceleración del bloque.
}